\section{Which singularities can persist?}\label{sec:strata}

Choose one of the two local components at each degree-six vertex and
denote this choice by $S$. Let $V_S$ be the direct sum of the chosen
summands and the full spaces at all other points. Put $K_S=K\cap V_S$.
On $V_S$ take the following linear forms, with a separate index for
each marked cycle. At an original node take its coordinate. At a
degree-six vertex take the two or three signed nodal periods $q_j$ of
the chosen modification. At a degree-five vertex take the five forms
$\lambda_i$ of Lemma~\ref{lem:linkperiods}. Denote the resulting list
by $(\ell_j)_{j\in J_S}$. Its support at a point is the subset on
which the evaluations are nonzero.

\begin{corollary}\label{cor:strata}
The reduction of $\Def(X)$ is the product of $\C^h$ with the reduced
union of the $K_S$. The scheme-theoretic intersection of any nonempty collection
of selected bases equals $\C^h$ times the intersection of the corresponding
$K_S$. For $I\subset J_S$ put
\[
 F_I=K_S\cap\bigcap_{j\notin I}\ker\ell_j .
\]
The support $I$ occurs precisely when $\ell_i|_{F_I}\ne0$ for every
$i\in I$. Its stratum has dimension $h+\dim F_I$ and closure
$\C^h\times F_I$.
\end{corollary}
\begin{proof}
On taking the reduction of the intersection in Theorem~\ref{thm:main},
its support is the union of the linear sections $K_S$. The sum of the
ideals of any selected components is generated on $K$ by their defining
linear forms, since these already imply the local quadratic equations.
This gives the asserted intersection with its scheme structure.
Inside $F_I$ the required support is the complement of the kernels
of $\ell_i$ for $i\in I$. A complex vector space is never the union
of finitely many proper linear subspaces. This proves both existence and
density, including for empty $I$.
\end{proof}

Within a choice $S$ the closure contains exactly the feasible
subsets of $I$. Where two choices differ, their local link spaces meet
at zero, which represents the original cone. The preceding intersection
formula therefore also identifies strata represented on several choices.
A zero coordinate at an original node retains that node. At a nonzero
parameter for a vertex there is one residual node per vanishing marked
form, while the zero block retains the vertex itself.
Every feasible support is attained on a disc: take a line in $F_I$
through a point of that support and use Theorem~\ref{thm:main}.
The disc can be chosen projective over its base since an ample line
bundle on $X$ extends locally in the deformation when $H^2(X,\OO_X)=0$.

When only nodes and degree-six vertices occur, let $C_j$ be the
marked curves of the original and incidence nodes on the chosen
modification, represented on $Y$ as in Lemma~\ref{lem:linkperiods}.
In the coefficient coordinates $q_j$, $K_S$ is the kernel of
$(q_j)\mapsto\sum q_j[C_j]$ in $H_2(Y;\C)$. Indeed pairing with
each exceptional surface imposes the local sum-zero condition, and
then the remaining equation is \eqref{eq:kernel}. Thus, writing
$r_I=\dim\langle[C_i]:i\in I\rangle$, the dimension is
$h+|I|-r_I$. The existence criterion says that removing any single
$i\in I$ leaves $r_I$ unchanged. For a nodal threefold and $I$ the
set of all nodes, this is precisely the existence of a relation with
all coefficients nonzero in Friedman's criterion.
